% GENERATED FILE. Edit canonical lessons, then run npm run generate:books.
% canonical-source-hash: fnv1a64:363ccda845584819
% canonical-lessons: PA-C313-viah, PA-C313-barat, PA-C313-lara, PA-C313-vandna, PA-C313-vadhai

\chapter{A Wedding in the Family}
\label{ch:pa-a2-a-wedding-in-the-family}
\hlchaptermodality{313}

\begin{chapteropening}
\textbf{By the end of this chapter:} \emph{I can say what happened at a family wedding.}

Complete the last lesson of chapter 313: I can say what happened at a family wedding.
\end{chapteropening}

\section[\texorpdfstring{\pa{ਵਿਆਹ} (viāh)}{ਵਿਆਹ (viāh)}]{\pa{ਵਿਆਹ} (viāh) — a wedding}

\label{lesson:PA-C313-viah}



\begin{quote}
Before the new one: say the Punjabi for a host, then the Punjabi for a gift.
\end{quote}

\subsection*{You'll want to know: \pa{ਵਿਆਹ}}
\textbf{\pa{ਵਿਆਹ}} — \emph{viāh} — \textquotedblleft{}a wedding\textquotedblright{}.

\begin{quote}
\pa{ਪਿਛਲੇ} \pa{ਮਹੀਨੇ} \pa{ਮੇਰੇ} \pa{ਭਰਾ} \pa{ਦਾ} \pa{ਵਿਆਹ} \pa{ਸੀ।}
\end{quote}

\begin{quote}
\emph{pichhle mahīne mere bharā dā viāh sī.}
\end{quote}

\begin{quote}
My brother's wedding was last month.
\end{quote}

A masculine noun.

\begin{cousinweb}
\textbf{\pa{ਵਿਆਹ}} (\emph{viāh}) is inherited from Sanskrit \emph{vivāha}, \textquotedblleft{}marriage\textquotedblright{}; CDIAL lists Punjabi \emph{viāh}, \textquotedblleft{}marriage\textquotedblright{}, under it (CDIAL 11920).
\end{cousinweb}

\begin{grammarlens}[title={What you've built}]
One more word for saying what happened.
\end{grammarlens}

\subsection*{Your turn}
\begin{itemize}
\raggedright
  \item \emph{Say it:} \emph{viāh}
  \item \emph{Say it:} \emph{viāh}, once more
  \item \emph{Say it:} \emph{pichhle mahīne mere bharā dā viāh sī.} — My brother's wedding was last month.
  \item \emph{Recall:} say the Punjabi for a host, then the Punjabi for a gift, then say \emph{viāh} again
\end{itemize}

\begin{tcolorbox}[breakable,colback=teal!4,colframe=teal!35!black,title={Before you move on}]
What does \pa{ਵਿਆਹ} mean? (\textquotedblleft{}A wedding\textquotedblright{}.) Say it once more.
\end{tcolorbox}
\section[\texorpdfstring{\pa{ਬਰਾਤ} (barāt)}{ਬਰਾਤ (barāt)}]{\pa{ਬਰਾਤ} (barāt) — a groom's wedding party}

\label{lesson:PA-C313-barat}



\begin{quote}
Before the new one: say the Punjabi for a gift, then the Punjabi for a wedding.
\end{quote}

\subsection*{You'll want to know: \pa{ਬਰਾਤ}}
\textbf{\pa{ਬਰਾਤ}} — \emph{barāt} — \textquotedblleft{}a groom's wedding party\textquotedblright{}.

\begin{quote}
\pa{ਬਰਾਤ} \pa{ਸ਼ਾਮ} \pa{ਨੂੰ} \pa{ਆਈ} \pa{ਅਤੇ} \pa{ਸਭ} \pa{ਨੱਚੇ।}
\end{quote}

\begin{quote}
\emph{barāt shām nū̃ āī ate sabh nachche.}
\end{quote}

\begin{quote}
The groom's party arrived in the evening, and everyone danced.
\end{quote}

A feminine noun.

\begin{cousinweb}
\textbf{\pa{ਬਰਾਤ}} (\emph{barāt}) is inherited from Sanskrit \emph{varayātrā}, \textquotedblleft{}the procession of a suitor or bridegroom\textquotedblright{}; CDIAL lists Punjabi \emph{barāt}, \textquotedblleft{}groom's procession\textquotedblright{}, under it (CDIAL 11322).
\end{cousinweb}

\begin{grammarlens}[title={What you've built}]
One more word for saying what happened.
\end{grammarlens}

\subsection*{Your turn}
\begin{itemize}
\raggedright
  \item \emph{Say it:} \emph{barāt}
  \item \emph{Say it:} \emph{barāt}, once more
  \item \emph{Say it:} \emph{barāt shām nū̃ āī ate sabh nachche.} — The groom's party arrived in the evening, and everyone danced.
  \item \emph{Recall:} say the Punjabi for a gift, then the Punjabi for a wedding, then say \emph{barāt} again
  \item \emph{Write it:} copy \textbf{\pa{ਬਰਾਤ}} once, checking the whole word
\end{itemize}

\begin{tcolorbox}[breakable,colback=teal!4,colframe=teal!35!black,title={Before you move on}]
What does \pa{ਬਰਾਤ} mean? (\textquotedblleft{}A groom's wedding party\textquotedblright{}.) Say it once more.
\end{tcolorbox}
\section[\texorpdfstring{\pa{ਲਾੜਾ} (lāṛā)}{ਲਾੜਾ (lāṛā)}]{\pa{ਲਾੜਾ} (lāṛā) — a bridegroom}

\label{lesson:PA-C313-lara}



\begin{quote}
Before the new one: say the Punjabi for a wedding, then the Punjabi for a groom's wedding party.
\end{quote}

\subsection*{You'll want to know: \pa{ਲਾੜਾ}}
\textbf{\pa{ਲਾੜਾ}} — \emph{lāṛā} — \textquotedblleft{}a bridegroom\textquotedblright{}.

\begin{quote}
\pa{ਲਾੜਾ} \pa{ਬਹੁਤ} \pa{ਖ਼ੁਸ਼} \pa{ਸੀ} \pa{ਅਤੇ} \pa{ਸਭ} \pa{ਨਾਲ} \pa{ਨੱਚਿਆ।}
\end{quote}

\begin{quote}
\emph{lāṛā bahut khush sī ate sabh nāl nachchiā.}
\end{quote}

\begin{quote}
The bridegroom was very happy and danced with everyone.
\end{quote}

A masculine noun.

\begin{cousinweb}
\textbf{\pa{ਲਾੜਾ}} (\emph{lāṛā}) is inherited from a reconstructed \emph{*lāḍa}, \textquotedblleft{}fondling\textquotedblright{}; CDIAL lists Punjabi \emph{lāṛā}, \textquotedblleft{}bridegroom\textquotedblright{}, under it (CDIAL 11012).
\end{cousinweb}

\begin{grammarlens}[title={What you've built}]
One more word for saying what happened.
\end{grammarlens}

\subsection*{Your turn}
\begin{itemize}
\raggedright
  \item \emph{Say it:} \emph{lāṛā}
  \item \emph{Say it:} \emph{lāṛā}, once more
  \item \emph{Say it:} \emph{lāṛā bahut khush sī ate sabh nāl nachchiā.} — The bridegroom was very happy and danced with everyone.
  \item \emph{Recall:} say the Punjabi for a wedding, then the Punjabi for a groom's wedding party, then say \emph{lāṛā} again
  \item \emph{Write it:} copy \textbf{\pa{ਲਾੜਾ}} once, checking the whole word
\end{itemize}

\begin{tcolorbox}[breakable,colback=teal!4,colframe=teal!35!black,title={Before you move on}]
What does \pa{ਲਾੜਾ} mean? (\textquotedblleft{}A bridegroom\textquotedblright{}.) Say it once more.
\end{tcolorbox}
\section[\texorpdfstring{\pa{ਵੰਡਣਾ} (vaṇḍṇā)}{ਵੰਡਣਾ (vaṇḍṇā)}]{\pa{ਵੰਡਣਾ} (vaṇḍṇā) — to hand out}

\label{lesson:PA-C313-vandna}



\begin{quote}
Before the new one: say the Punjabi for a groom's wedding party, then the Punjabi for a bridegroom.
\end{quote}

\subsection*{You'll want to know: \pa{ਵੰਡਣਾ}}
\textbf{\pa{ਵੰਡਣਾ}} — \emph{vaṇḍṇā} — \textquotedblleft{}to hand out\textquotedblright{}.

\begin{quote}
\pa{ਵਿਆਹ} \pa{ਵਿੱਚ} \pa{ਅਸੀਂ} \pa{ਸਭ} \pa{ਨੂੰ} \pa{ਖੀਰ} \pa{ਵੰਡੀ।}
\end{quote}

\begin{quote}
\emph{viāh vich asīṁ sabh nū̃ khīr vaṇḍī.}
\end{quote}

\begin{quote}
At the wedding we handed out rice pudding to everyone.
\end{quote}

Past (I …): \pa{ਮੈਂ} \pa{ਦੋਸਤਾਂ} \pa{ਨੂੰ} \pa{ਕੇਲੇ} \pa{ਵੰਡੇ}, \emph{maiṁ dostāṁ nū̃ kele vaṇḍe}.

\begin{cousinweb}
\textbf{\pa{ਵੰਡਣਾ}} (\emph{vaṇḍṇā}) is inherited from Sanskrit \emph{vaṇṭati}, \textquotedblleft{}shares\textquotedblright{}; CDIAL lists Punjabi \emph{vaṇḍṇā}, \textquotedblleft{}shares\textquotedblright{}, under it (CDIAL 11238).
\end{cousinweb}

\begin{grammarlens}[title={What you've built}]
One more word for saying what happened.
\end{grammarlens}

\subsection*{Your turn}
\begin{itemize}
\raggedright
  \item \emph{Say it:} \emph{vaṇḍṇā}
  \item \emph{Say it:} \emph{vaṇḍṇā}, once more
  \item \emph{Say it:} \emph{viāh vich asīṁ sabh nū̃ khīr vaṇḍī.} — At the wedding we handed out rice pudding to everyone.
  \item \emph{Recall:} say the Punjabi for a groom's wedding party, then the Punjabi for a bridegroom, then say \emph{vaṇḍṇā} again
  \item \emph{Write it:} copy \textbf{\pa{ਵੰਡਣਾ}} once, checking the whole word
\end{itemize}

\begin{tcolorbox}[breakable,colback=teal!4,colframe=teal!35!black,title={Before you move on}]
What does \pa{ਵੰਡਣਾ} mean? (\textquotedblleft{}To hand out\textquotedblright{}.) Say it once more.
\end{tcolorbox}
\section[\texorpdfstring{\pa{ਵਧਾਈ} (vadhāī)}{ਵਧਾਈ (vadhāī)}]{\pa{ਵਧਾਈ} (vadhāī) — congratulations}

\label{lesson:PA-C313-vadhai}



\begin{quote}
Before the new one: say the Punjabi for a bridegroom, then the Punjabi for to hand out.
\end{quote}

\subsection*{You'll want to know: \pa{ਵਧਾਈ}}
\textbf{\pa{ਵਧਾਈ}} — \emph{vadhāī} — \textquotedblleft{}congratulations\textquotedblright{}.

\begin{quote}
\pa{ਮੈਂ} \pa{ਲਾੜੇ} \pa{ਨੂੰ} \pa{ਵਧਾਈ} \pa{ਦਿੱਤੀ।}
\end{quote}

\begin{quote}
\emph{maiṁ lāṛe nū̃ vadhāī dittī.}
\end{quote}

\begin{quote}
I congratulated the bridegroom.
\end{quote}

A feminine noun.

\begin{cousinweb}
\textbf{\pa{ਵਧਾਈ}} (\emph{vadhāī}) is inherited from Sanskrit \emph{vardhāpayati}, \textquotedblleft{}increases, gladdens, congratulates\textquotedblright{}; CDIAL lists Punjabi \emph{vadhāī}, \textquotedblleft{}blessing\textquotedblright{}, under it (CDIAL 11383).
\end{cousinweb}

\begin{grammarlens}[title={What you've built}]
One more word for saying what happened.
\end{grammarlens}

\subsection*{Your turn}
\begin{itemize}
\raggedright
  \item \emph{Say it:} \emph{vadhāī}
  \item \emph{Say it:} \emph{vadhāī}, once more
  \item \emph{Say it:} \emph{maiṁ lāṛe nū̃ vadhāī dittī.} — I congratulated the bridegroom.
  \item \emph{Recall:} say the Punjabi for a bridegroom, then the Punjabi for to hand out, then say \emph{vadhāī} again
  \item \emph{Write it:} copy \textbf{\pa{ਵਧਾਈ}} once, checking the whole word
\end{itemize}

\begin{tcolorbox}[breakable,colback=teal!4,colframe=teal!35!black,title={Before you move on}]
What does \pa{ਵਧਾਈ} mean? (\textquotedblleft{}Congratulations\textquotedblright{}.) Say it once more.
\end{tcolorbox}
